\documentclass[10pt,a4paper]{amsart}
\usepackage[margin=19mm]{geometry}
\usepackage{amsmath,amssymb,mathtools}
\usepackage[scr=rsfs,cal=euler]{mathalfa}
\usepackage{xfrac}
\usepackage[unicode,hidelinks,bookmarks=false]{hyperref}
\newcommand{\ie}{i.e.\ }
\newcommand{\cf}{cf.\ }
\newcommand{\resp}{respectively\ }
\newcommand{\Subsection}{\subsection}
\providecommand{\nobreakdash}{\nobreak\hbox{-}}
\setlength{\emergencystretch}{2em}
\multlinegap0pt
\usepackage{xcolor}
\usepackage{amssymb}
\usepackage{enumerate}
\newtheorem{theorem}{Theorem}
\title{On matrix valued (asymmetric) truncated Toeplitz operators}
\author{NİHAT GÖKHAN GÖĞÜŞ, REWAYAT KHAN}
\date{Source: arXiv:2603.23579v1 (CC BY 4.0)}
\begin{document}
\maketitle

\noindent\emph{Source and licence.} On matrix valued (asymmetric) truncated Toeplitz operators. Version arXiv:2603.23579v1 (CC BY 4.0), distributed by arXiv under the Creative Commons Attribution 4.0 International licence, http://creativecommons.org/licenses/by/4.0/, as recorded in the arXiv licence metadata for this exact version and shown on the abstract page https://arxiv.org/abs/2603.23579v1. Full licence notice: \texttt{licenses/arxiv-2603.23579-CC-BY-4.0.txt}.\par\smallskip

\noindent\textit{Editorial scope.} Counted unit: the article's theorem conj (source0.tex lines 400--408). The article's complete original proof of it is delivered with it (source0.tex lines 409--443, one \texttt{proof} environment of 35 lines). No mathematics is written, repaired or re-derived: the statement and the proof are the article's own lines, in the article's own notation, and no retained line is substituted or re-flowed.  No further source-local context is needed: every cross-reference of every retained line resolves inside the two delivered spans.  Nothing was omitted from any retained span, so the package carries no omission record.  No retained line contains a citation, so the wrapper carries no bibliography and no bibliography entry.  No retained line contains a figure environment, an image-inclusion command or drawing code, so no image is delivered and the package declares no asset dependency.  Editorial formatting only, in addition: an AMS \texttt{amsart} preamble that restates the article's own declarations and macros used by the retained lines, namely the environments \texttt{theorem} on separate counters, together with the standard packages those lines need.  The retained environments print in this document as Theorem 1 in order of appearance, whereas the article prints the same items with the numbering of its own full document.  The complete native source of the article as distributed in the arXiv e-print for this exact version is bundled unchanged as \texttt{sources/197091/source0.tex}.
\begin{theorem}\label{conj}
Let $\Theta, \Lambda$, and $\Psi$ be $\Gamma$-symmetric inner functions such that $\Theta=\Lambda\Psi$. Then, the operator $$C_{\Lambda,\Psi}:K_{\Theta}=K_{\Lambda}\oplus\Lambda K_{\Psi}\to K_{\Theta}=K_{\Lambda}\oplus\Lambda K_{\Psi}$$ defined by
$$C_{\Lambda,\Psi}=C_{\Lambda}\oplus\Lambda C_{\Psi}\Lambda^{*}=\begin{bmatrix}
C_{\Lambda}&0\\
0&\Lambda C_{\Psi}\Lambda^{*}
\end{bmatrix}
$$
is a conjugation on $K_{\Theta}$.
\end{theorem}

\begin{proof}
To show that $C_{\Lambda,\Psi}$ is a conjugation on $K_{\Theta}$, we have to show that $\Lambda C_{\Psi}\Lambda^*$ is a conjugation on $\Lambda K_{\Psi}$. Since $\Lambda$ and $\Psi$ are $\Gamma$-symmetric, we have $C_{\Lambda}$ and $C_{\Psi}$ as conjugations on $K_{\Lambda}$ and $K_{\Psi}$, respectively. Let $\mathcal{C}=\Lambda C_{\Psi}\Lambda^*$. First we have to show that $\mathcal{C}(\Lambda K_\Psi)=\Lambda K_\Psi$.
Let $f\in \Lambda K_\Psi$, then $f=\Lambda g$ for some $g\in K_\Psi$. Then
$$\mathcal{C}(f)=\Lambda C_\Psi\Lambda^*(\Lambda g)=\Lambda C_\Psi g \in \Lambda K_\Psi.$$ It follows that $\mathcal{C}(\Lambda K_\Psi) \subset  \Lambda K_\Psi$.
The operator $\mathcal{C}$ is involution because $$\mathcal{C}^2=\Lambda C_{\Psi}\Lambda^* \Lambda C_{\Psi}\Lambda^*=\Lambda C_{\Psi}^2\Lambda^*=I_{\Lambda K_\Psi}.$$
Hence, we have
$$\mathcal{C}(\Lambda K_\Psi)=\Lambda K_\Psi.$$
The antilinearity of $\mathcal{C}$ follows from the antilinearity of $C_\Psi$.

Now, we prove that $\mathcal{C}$ is isometric, i.e., $\|\mathcal{C}f\|=\|f\|$ for all $f=\Lambda g\in  \Lambda K_{\Psi}$, where $g\in K_{\Psi}$. Consider

\begin{align*}
\|\mathcal{C}f\|^{2}
&=\langle
\mathcal{C}f, \mathcal{C}f
\rangle\\
&=\langle
\Lambda C_\Psi \Lambda^*\Lambda g, \Lambda C_\Psi \Lambda^*\Lambda g
\rangle \\
&=\langle
 C_\Psi g, C_\Psi g
\rangle \\
&=\langle
 g, g
\rangle\\
&=\langle
 \Lambda g, \Lambda g
\rangle\\
&=\langle
 f, f
\rangle\\
&=\|f\|^{2}.
\end{align*}
Hence, $\mathcal{C}$ is a conjugation, which follows that $C_{\Lambda, \Psi}$ is a conjugation.
\end{proof}

\end{document}
